\begin{PSbox}[unbreakable]{obj}{اهداف یادگیری}

پس از تکمیل این ماژول، دانشجویان قادر خواهند بود:

\begin{enumerate}
\def\labelenumi{\arabic{enumi}.}
\tightlist
\item
  قانون اعداد بزرگ را بیان و توضیح دهند
\item
  قضیه حد مرکزی را بیان و به کار ببرند
\item
  توضیح دهند چرا توزیع‌های گاوسی این‌قدر در مهندسی فراوان‌اند
\item
  از \lr{\mbox{CLT}} برای تقریب‌زنی استفاده کنند
\item
  \lr{\mbox{CLT}} را به‌صورت بصری با شبیه‌سازی‌های مونت‌کارلوی \lr{\mbox{MATLAB}} نشان دهند
\end{enumerate}

\end{PSbox}

\PSsection{11.1}{مقدمه: چرا گاوسی همه‌جا هست؟}

توزیع گاوسی در موارد زیر ظاهر می‌شود:

\begin{itemize}
\tightlist
\item
  نویز حرارتی، خطاهای اندازه‌گیری، رواداری‌های ساخت، تداخل تجمعی، نوسان‌های سهام
\end{itemize}

\textbf{دلیل:} همه این کمیت‌ها جمع‌ها (یا میانگین‌ها)ی بسیاری از اثرهای تصادفی کوچک و مستقل هستند. \lr{\mbox{CLT}} همگرایی به گاوسی را صرف‌نظر از توزیع زیربنایی تضمین می‌کند.

\PSsection{11.2}{قانون اعداد بزرگ \lr{\mbox{(LLN)}}}

\PSsubsection{قانون ضعیف اعداد بزرگ}

برای \lr{$X_{1},\, X_{2},\, \ldots,\, X_{n}$} مستقل و هم‌توزیع \lr{$(i.i.d.)$} با میانگین \lr{$\mu$} و واریانس متناهی \lr{$\sigma^{2}$}:

\PSdisplay{\bar{X}_n = \frac{1}{n}\sum_{i=1}^n X_i \xrightarrow{P} \mu \quad \PSmtext{as } n \to \infty}

(«همگرایی در احتمال»: \lr{$P(\lvert \bar{X}_{n} - \mu\rvert > \varepsilon) \to 0$} برای هر \lr{$\varepsilon > 0$})

\PSsubsection{قانون قوی اعداد بزرگ}

\PSdisplay{\bar{X}_n \xrightarrow{a.s.} \mu \quad \PSmtext{as } n \to \infty}

(«همگرایی تقریباً مطمین»: \lr{$P(\bar{X}_{n} \to \mu) = 1$})

\PSsubsection{اهمیت مهندسی}

\begin{itemize}
\tightlist
\item
  \textbf{میانگین نمونه با رشد اندازه نمونه به میانگین واقعی همگرا می‌شود}
\item
  برآورد میانگین‌ها از داده را توجیه می‌کند
\item
  بنیاد شبیه‌سازی مونت‌کارلو (شبیه‌سازی بارها \lr{\mbox{\ensuremath{\to}}} میانگین همگرا می‌شود)
\item
  این گزاره را توجیه می‌کند: «اگر به‌قدر کافی اندازه‌گیری کنم، میانگین من به مقدار واقعی نزدیک خواهد بود»
\end{itemize}

\PSsubsection{با چه سرعتی؟}

بر اساس چبیشف: \lr{$P(\lvert \bar{X}_{n} - \mu\rvert > \varepsilon) \le \frac{\sigma^{2}}{n\varepsilon^{2}}$}

برای اطمینان 95\% از این‌که خطا \lr{$< \varepsilon$} باشد: به \lr{$n \ge \frac{\sigma^{2}}{0.05 \cdot \varepsilon^{2}}$} نیاز است

\PSsection{11.3}{قضیه حد مرکزی \lr{\mbox{(CLT)}}}

\PSsubsection{بیان صوری}

برای \lr{$X_{1},\, X_{2},\, \ldots,\, X_{n}$} مستقل و هم‌توزیع \lr{$(i.i.d.)$} با میانگین \lr{$\mu$} و واریانس متناهی \lr{$\sigma^{2}$}:

\PSdisplay{Z_n = \frac{\bar{X}_n - \mu}{\sigma/\sqrt{n}} = \frac{\sum_{i=1}^n X_i - n\mu}{\sigma\sqrt{n}} \xrightarrow{d} N(0,1)}

\PSsubsection{شکل عملی}

برای \lr{$n$} بزرگ:\PSdisplay{\bar{X}_n \approx N\left(\mu, \frac{\sigma^2}{n}\right)}

\PSdisplay{S_n = \sum_{i=1}^n X_i \approx N(n\mu, n\sigma^2)}

\PSsubsection{شرایط}

\begin{enumerate}
\def\labelenumi{\arabic{enumi}.}
\tightlist
\item
  \lr{$X_{1},\, \ldots,\, X_{n}$} مستقل هستند
\item
  هم‌توزیع هستند (با \lr{\mbox{CLT}} لیپانوف/لینده‌برگ می‌توان این شرط را سست کرد)
\item
  میانگین متناهی \lr{$\mu$} و واریانس متناهی \lr{$\sigma^{2}$}
\item
  \lr{$n$} «به‌قدر کافی بزرگ» است (قاعده سرانگشتی: \lr{$n \ge 30$}، اما به شکل توزیع بستگی دارد)
\end{enumerate}

\PSsubsection{آنچه \lr{\mbox{CLT}} نیاز ندارد}

\begin{itemize}
\tightlist
\item
  نیاز ندارد \lr{$X_{i}$} گاوسی باشد
\item
  نیاز ندارد \lr{$X_{i}$} پیوسته باشد
\item
  نیاز ندارد هیچ توزیع خاصی داشته باشد
\end{itemize}

\PSsection{11.4}{نرخ همگرایی}

\lr{\mbox{CLT}} با چه سرعتی «اثر می‌گذارد»؟

{\def\LTcaptype{none} % do not increment counter
\begin{longtable}[c]{@{}cc@{}}
\toprule\noalign{}
\rowcolor{PSnavy}\PSth{توزیع اصلی} & \PSth{\lr{$n$} لازم برای تقریب خوب} \\
\midrule\noalign{}
\endhead
\bottomrule\noalign{}
\endlastfoot
متقارن (یکنواخت) & \lr{$n \approx 5-10$} \\
کمی چوله & \lr{$n \approx 20-30$} \\
به‌شدت چوله (نمایی) & \lr{$n \approx 30-50$} \\
به‌غایت چوله & \lr{$n \approx 50-100+$} \\
\end{longtable}
}

\textbf{قاعده:} توزیع‌های متقارن‌تر \lr{\mbox{\ensuremath{\to}}} همگرایی سریع‌تر.

\PSsection{11.5}{نمایش بصری}

\PSsubsection{جمع متغیرهای تصادفی یکنواخت \lr{\mbox{\ensuremath{\to}}} گاوسی}

\begin{itemize}
\tightlist
\item
  \lr{$n = 1$}: یکنواخت (مسطح)
\item
  \lr{$n = 2$}: مثلثی
\item
  \lr{$n = 3$}: شروع به زنگوله‌ای‌شدن می‌کند
\item
  \lr{$n = 5$}: تقریباً گاوسی
\item
  \lr{$n = 12$}: عملاً از گاوسی قابل تفکیک نیست
\end{itemize}

\PSsubsection{جمع متغیرهای تصادفی نمایی \lr{\mbox{\ensuremath{\to}}} گاوسی}

\begin{itemize}
\tightlist
\item
  \lr{$n = 1$}: نمایی (به‌شدت چوله به راست)
\item
  \lr{$n = 3$}: همچنان چوله اما کمتر \lr{$(\mathrm{Gamma}(3))$}
\item
  \lr{$n = 10$}: تقریباً متقارن، زنگوله‌ای
\item
  \lr{$n = 30$}: اساساً گاوسی
\end{itemize}

\PSsubsection{دوجمله‌ای \lr{\mbox{\ensuremath{\to}}} تقریب نرمال}

\lr{$\mathrm{Binomial}(n,\, p) \approx N(np,\, np(1-p))$} برای \lr{$n$} بزرگ

این همان \lr{\mbox{CLT}} اعمال‌شده به جمع آزمایش‌های برنولی است!

\begin{itemize}
\tightlist
\item
  تصحیح پیوستگی: \lr{$P(X \le k) \approx \Phi \left(\frac{k + 0.5 - np}{\sqrt{np(1-p)}}\right)$}
\end{itemize}

\PSsection{11.6}{اهمیت مهندسی}

\PSsubsection{چرا نویز حرارتی گاوسی است}

نویز حرارتی = اثر تجمعی میلیاردها الکترون که به‌طور تصادفی حرکت می‌کنند. سهم هر الکترون ناچیز و مستقل است \lr{\mbox{\ensuremath{\to} CLT}} اعمال می‌شود \lr{\mbox{\ensuremath{\to}}} نویز گاوسی است.

\PSsubsection{چرا تداخل تجمعی گاوسی است}

در یک شبکه سلولی، تداخل کل = جمع سیگنال‌های بسیاری از کاربران مستقل. \lr{\mbox{$\mathrm{CLT}$ \ensuremath{\to}}} تداخل \lr{\mbox{\ensuremath{\approx}}} گاوسی (برای کاربران زیاد).

\PSsubsection{چرا خطاهای اندازه‌گیری گاوسی هستند}

خطای کل = جمع بسیاری از منابع خطای کوچک و مستقل (کالیبراسیون، کوانتیزاسیون، گرمایی، لرزش و \lr{…}) \lr{\mbox{$\to \mathrm{CLT}$ \ensuremath{\to}}} خطا \lr{\mbox{\ensuremath{\approx}}} گاوسی.

\PSsubsection{چرا تغییرات ساخت گاوسی هستند}

مقدار مقاومت = مقدار نامی + جمع بسیاری از تغییرات کوچک فرایند (دوپینگ، حکاکی، دما و \lr{…}) \lr{\mbox{$\to \mathrm{CLT}$ \ensuremath{\to}}} مقدار \lr{\mbox{\ensuremath{\approx}}} گاوسی حول مقدار نامی.

\PSsection{11.7}{نمایش‌های مونت‌کارلوی \lr{\mbox{MATLAB}}}

\begin{PSbox}[unbreakable]{ex}{مثال 1: \lr{\mbox{CLT}} با توزیع یکنواخت}

\tcbset{PScodemode/.style={unbreakable}}

\begin{Shaded}
\begin{Highlighting}[]
\CommentTok{\%\% CLT: Sum of Uniform RVs approaches Gaussian}
\VariableTok{N\_experiments} \OperatorTok{=} \FloatTok{50000}\OperatorTok{;}
\VariableTok{n\_values} \OperatorTok{=}\NormalTok{ [}\FloatTok{1}\OperatorTok{,} \FloatTok{2}\OperatorTok{,} \FloatTok{5}\OperatorTok{,} \FloatTok{12}\OperatorTok{,} \FloatTok{30}\NormalTok{]}\OperatorTok{;}

\VariableTok{figure}\OperatorTok{;}
\KeywordTok{for} \VariableTok{idx} \OperatorTok{=} \FloatTok{1}\OperatorTok{:}\VariableTok{length}\NormalTok{(}\VariableTok{n\_values}\NormalTok{)}
    \VariableTok{n} \OperatorTok{=} \VariableTok{n\_values}\NormalTok{(}\VariableTok{idx}\NormalTok{)}\OperatorTok{;}
    
    \CommentTok{\% Generate sum of n Uniform(0,1) RVs}
    \VariableTok{X} \OperatorTok{=} \VariableTok{sum}\NormalTok{(}\VariableTok{rand}\NormalTok{(}\VariableTok{n}\OperatorTok{,} \VariableTok{N\_experiments}\NormalTok{)}\OperatorTok{,} \FloatTok{1}\NormalTok{)}\OperatorTok{;}
    
    \CommentTok{\% Standardize: Z = (X {-} nμ)/(σ√n) where μ=0.5, σ²=1/12}
    \VariableTok{mu\_sum} \OperatorTok{=} \VariableTok{n} \OperatorTok{*} \FloatTok{0.5}\OperatorTok{;}
    \VariableTok{sigma\_sum} \OperatorTok{=} \VariableTok{sqrt}\NormalTok{(}\VariableTok{n} \OperatorTok{/} \FloatTok{12}\NormalTok{)}\OperatorTok{;}
    \VariableTok{Z} \OperatorTok{=}\NormalTok{ (}\VariableTok{X} \OperatorTok{{-}} \VariableTok{mu\_sum}\NormalTok{) }\OperatorTok{/} \VariableTok{sigma\_sum}\OperatorTok{;}
    
    \VariableTok{subplot}\NormalTok{(}\FloatTok{2}\OperatorTok{,} \FloatTok{3}\OperatorTok{,} \VariableTok{idx}\NormalTok{)}\OperatorTok{;}
    \VariableTok{histogram}\NormalTok{(}\VariableTok{Z}\OperatorTok{,} \FloatTok{60}\OperatorTok{,} \SpecialStringTok{\textquotesingle{}Normalization\textquotesingle{}}\OperatorTok{,} \SpecialStringTok{\textquotesingle{}pdf\textquotesingle{}}\NormalTok{)}\OperatorTok{;}
    \VariableTok{hold} \VariableTok{on}\OperatorTok{;}
    \VariableTok{z} \OperatorTok{=} \VariableTok{linspace}\NormalTok{(}\OperatorTok{{-}}\FloatTok{4}\OperatorTok{,} \FloatTok{4}\OperatorTok{,} \FloatTok{200}\NormalTok{)}\OperatorTok{;}
    \VariableTok{plot}\NormalTok{(}\VariableTok{z}\OperatorTok{,} \VariableTok{normpdf}\NormalTok{(}\VariableTok{z}\NormalTok{)}\OperatorTok{,} \SpecialStringTok{\textquotesingle{}r\textquotesingle{}}\OperatorTok{,} \SpecialStringTok{\textquotesingle{}LineWidth\textquotesingle{}}\OperatorTok{,} \FloatTok{2}\NormalTok{)}\OperatorTok{;}
    \VariableTok{title}\NormalTok{(}\VariableTok{sprintf}\NormalTok{(}\SpecialStringTok{\textquotesingle{}n = \%d\textquotesingle{}}\OperatorTok{,} \VariableTok{n}\NormalTok{))}\OperatorTok{;}
    \VariableTok{xlabel}\NormalTok{(}\SpecialStringTok{\textquotesingle{}Standardized Sum\textquotesingle{}}\NormalTok{)}\OperatorTok{;} \VariableTok{xlim}\NormalTok{([}\OperatorTok{{-}}\FloatTok{4} \FloatTok{4}\NormalTok{])}\OperatorTok{;}
    \KeywordTok{if} \VariableTok{idx} \OperatorTok{==} \FloatTok{1}\OperatorTok{,} \VariableTok{ylabel}\NormalTok{(}\SpecialStringTok{\textquotesingle{}PDF\textquotesingle{}}\NormalTok{)}\OperatorTok{;} \KeywordTok{end}
\KeywordTok{end}
\VariableTok{sgtitle}\NormalTok{(}\SpecialStringTok{\textquotesingle{}CLT: Sum of Uniform RVs → Gaussian\textquotesingle{}}\NormalTok{)}\OperatorTok{;}
\end{Highlighting}
\end{Shaded}

\end{PSbox}

\begin{PSbox}[unbreakable]{ex}{مثال 2: \lr{\mbox{CLT}} با توزیع نمایی}

\tcbset{PScodemode/.style={unbreakable}}

\begin{Shaded}
\begin{Highlighting}[]
\CommentTok{\%\% CLT: Sample Mean of Exponential(1) converges to Gaussian}
\VariableTok{lambda} \OperatorTok{=} \FloatTok{1}\OperatorTok{;}  \VariableTok{mu} \OperatorTok{=} \FloatTok{1}\OperatorTok{;}  \VariableTok{sigma} \OperatorTok{=} \FloatTok{1}\OperatorTok{;}
\VariableTok{N\_experiments} \OperatorTok{=} \FloatTok{100000}\OperatorTok{;}
\VariableTok{n\_values} \OperatorTok{=}\NormalTok{ [}\FloatTok{1}\OperatorTok{,} \FloatTok{3}\OperatorTok{,} \FloatTok{10}\OperatorTok{,} \FloatTok{30}\OperatorTok{,} \FloatTok{100}\NormalTok{]}\OperatorTok{;}

\VariableTok{figure}\OperatorTok{;}
\KeywordTok{for} \VariableTok{idx} \OperatorTok{=} \FloatTok{1}\OperatorTok{:}\VariableTok{length}\NormalTok{(}\VariableTok{n\_values}\NormalTok{)}
    \VariableTok{n} \OperatorTok{=} \VariableTok{n\_values}\NormalTok{(}\VariableTok{idx}\NormalTok{)}\OperatorTok{;}
    
    \CommentTok{\% Generate sample means}
    \VariableTok{X} \OperatorTok{=} \VariableTok{exprnd}\NormalTok{(}\FloatTok{1}\OperatorTok{,} \VariableTok{n}\OperatorTok{,} \VariableTok{N\_experiments}\NormalTok{)}\OperatorTok{;}
    \VariableTok{X\_bar} \OperatorTok{=} \VariableTok{mean}\NormalTok{(}\VariableTok{X}\OperatorTok{,} \FloatTok{1}\NormalTok{)}\OperatorTok{;}
    
    \VariableTok{subplot}\NormalTok{(}\FloatTok{2}\OperatorTok{,} \FloatTok{3}\OperatorTok{,} \VariableTok{idx}\NormalTok{)}\OperatorTok{;}
    \VariableTok{histogram}\NormalTok{(}\VariableTok{X\_bar}\OperatorTok{,} \FloatTok{80}\OperatorTok{,} \SpecialStringTok{\textquotesingle{}Normalization\textquotesingle{}}\OperatorTok{,} \SpecialStringTok{\textquotesingle{}pdf\textquotesingle{}}\NormalTok{)}\OperatorTok{;}
    \VariableTok{hold} \VariableTok{on}\OperatorTok{;}
    \VariableTok{x} \OperatorTok{=} \VariableTok{linspace}\NormalTok{(}\FloatTok{0}\OperatorTok{,} \VariableTok{max}\NormalTok{(}\VariableTok{X\_bar}\NormalTok{)}\OperatorTok{,} \FloatTok{200}\NormalTok{)}\OperatorTok{;}
    \VariableTok{plot}\NormalTok{(}\VariableTok{x}\OperatorTok{,} \VariableTok{normpdf}\NormalTok{(}\VariableTok{x}\OperatorTok{,} \VariableTok{mu}\OperatorTok{,} \VariableTok{sigma}\OperatorTok{/}\VariableTok{sqrt}\NormalTok{(}\VariableTok{n}\NormalTok{))}\OperatorTok{,} \SpecialStringTok{\textquotesingle{}r\textquotesingle{}}\OperatorTok{,} \SpecialStringTok{\textquotesingle{}LineWidth\textquotesingle{}}\OperatorTok{,} \FloatTok{2}\NormalTok{)}\OperatorTok{;}
    \VariableTok{title}\NormalTok{(}\VariableTok{sprintf}\NormalTok{(}\SpecialStringTok{\textquotesingle{}n = \%d\textquotesingle{}}\OperatorTok{,} \VariableTok{n}\NormalTok{))}\OperatorTok{;}
    \VariableTok{xlabel}\NormalTok{(}\SpecialStringTok{\textquotesingle{}Sample Mean\textquotesingle{}}\NormalTok{)}\OperatorTok{;}
\KeywordTok{end}
\VariableTok{sgtitle}\NormalTok{(}\SpecialStringTok{\textquotesingle{}CLT: Mean of Exponential → Gaussian\textquotesingle{}}\NormalTok{)}\OperatorTok{;}
\end{Highlighting}
\end{Shaded}

\end{PSbox}

\begin{PSbox}[unbreakable]{ex}{مثال 3: تقریب نرمال دوجمله‌ای}

\tcbset{PScodemode/.style={unbreakable}}

\begin{Shaded}
\begin{Highlighting}[]
\CommentTok{\%\% Binomial(n,p) approximated by Gaussian}
\VariableTok{n} \OperatorTok{=} \FloatTok{50}\OperatorTok{;} \VariableTok{p} \OperatorTok{=} \FloatTok{0.3}\OperatorTok{;}
\VariableTok{k} \OperatorTok{=} \FloatTok{0}\OperatorTok{:}\VariableTok{n}\OperatorTok{;}
\VariableTok{pmf\_binom} \OperatorTok{=} \VariableTok{binopdf}\NormalTok{(}\VariableTok{k}\OperatorTok{,} \VariableTok{n}\OperatorTok{,} \VariableTok{p}\NormalTok{)}\OperatorTok{;}

\CommentTok{\% Normal approximation}
\VariableTok{mu\_approx} \OperatorTok{=} \VariableTok{n}\OperatorTok{*}\VariableTok{p}\OperatorTok{;}
\VariableTok{sigma\_approx} \OperatorTok{=} \VariableTok{sqrt}\NormalTok{(}\VariableTok{n}\OperatorTok{*}\VariableTok{p}\OperatorTok{*}\NormalTok{(}\FloatTok{1}\OperatorTok{{-}}\VariableTok{p}\NormalTok{))}\OperatorTok{;}
\VariableTok{pdf\_normal} \OperatorTok{=} \VariableTok{normpdf}\NormalTok{(}\VariableTok{k}\OperatorTok{,} \VariableTok{mu\_approx}\OperatorTok{,} \VariableTok{sigma\_approx}\NormalTok{)}\OperatorTok{;}

\VariableTok{figure}\OperatorTok{;}
\VariableTok{bar}\NormalTok{(}\VariableTok{k}\OperatorTok{,} \VariableTok{pmf\_binom}\OperatorTok{,} \SpecialStringTok{\textquotesingle{}FaceAlpha\textquotesingle{}}\OperatorTok{,} \FloatTok{0.5}\NormalTok{)}\OperatorTok{;}
\VariableTok{hold} \VariableTok{on}\OperatorTok{;}
\VariableTok{plot}\NormalTok{(}\VariableTok{k}\OperatorTok{,} \VariableTok{pdf\_normal}\OperatorTok{,} \SpecialStringTok{\textquotesingle{}r{-}\textquotesingle{}}\OperatorTok{,} \SpecialStringTok{\textquotesingle{}LineWidth\textquotesingle{}}\OperatorTok{,} \FloatTok{2}\NormalTok{)}\OperatorTok{;}
\VariableTok{xlabel}\NormalTok{(}\SpecialStringTok{\textquotesingle{}k\textquotesingle{}}\NormalTok{)}\OperatorTok{;} \VariableTok{ylabel}\NormalTok{(}\SpecialStringTok{\textquotesingle{}Probability\textquotesingle{}}\NormalTok{)}\OperatorTok{;}
\VariableTok{title}\NormalTok{(}\VariableTok{sprintf}\NormalTok{(}\SpecialStringTok{\textquotesingle{}Binomial(\%d, \%.1f) vs Normal(\%.1f, \%.2f)\textquotesingle{}}\OperatorTok{,} \VariableTok{n}\OperatorTok{,} \VariableTok{p}\OperatorTok{,} \VariableTok{mu\_approx}\OperatorTok{,} \VariableTok{sigma\_approx}\OperatorTok{\^{}}\FloatTok{2}\NormalTok{))}\OperatorTok{;}
\VariableTok{legend}\NormalTok{(}\SpecialStringTok{\textquotesingle{}Binomial (exact)\textquotesingle{}}\OperatorTok{,} \SpecialStringTok{\textquotesingle{}Normal approximation\textquotesingle{}}\NormalTok{)}\OperatorTok{;}
\end{Highlighting}
\end{Shaded}

\end{PSbox}

\PSsubsection{مثال 4: نمایش کامل \lr{\mbox{CLT}} با مونت‌کارلو}

\tcbset{PScodemode/.style={unbreakable}}

\begin{Shaded}
\begin{Highlighting}[]
\CommentTok{\%\% Full CLT demonstration: non{-}Gaussian → Gaussian via averaging}
\CommentTok{\% Generate N samples from a VERY non{-}Gaussian distribution}
\CommentTok{\% (mixture: 80\% from Exp(1), 20\% from Exp(0.1))}
\VariableTok{N\_experiments} \OperatorTok{=} \FloatTok{100000}\OperatorTok{;}
\VariableTok{n\_samples} \OperatorTok{=} \FloatTok{50}\OperatorTok{;}

\CommentTok{\% Generate raw data}
\VariableTok{raw} \OperatorTok{=} \VariableTok{zeros}\NormalTok{(}\VariableTok{n\_samples}\OperatorTok{,} \VariableTok{N\_experiments}\NormalTok{)}\OperatorTok{;}
\KeywordTok{for} \VariableTok{i} \OperatorTok{=} \FloatTok{1}\OperatorTok{:}\VariableTok{N\_experiments}
    \VariableTok{mask} \OperatorTok{=} \VariableTok{rand}\NormalTok{(}\VariableTok{n\_samples}\OperatorTok{,} \FloatTok{1}\NormalTok{) }\OperatorTok{\textless{}} \FloatTok{0.8}\OperatorTok{;}
    \VariableTok{raw}\NormalTok{(}\VariableTok{mask}\OperatorTok{,} \VariableTok{i}\NormalTok{) }\OperatorTok{=} \VariableTok{exprnd}\NormalTok{(}\FloatTok{1}\OperatorTok{,} \VariableTok{sum}\NormalTok{(}\VariableTok{mask}\NormalTok{)}\OperatorTok{,} \FloatTok{1}\NormalTok{)}\OperatorTok{;}
    \VariableTok{raw}\NormalTok{(}\OperatorTok{\textasciitilde{}}\VariableTok{mask}\OperatorTok{,} \VariableTok{i}\NormalTok{) }\OperatorTok{=} \VariableTok{exprnd}\NormalTok{(}\FloatTok{10}\OperatorTok{,} \VariableTok{sum}\NormalTok{(}\OperatorTok{\textasciitilde{}}\VariableTok{mask}\NormalTok{)}\OperatorTok{,} \FloatTok{1}\NormalTok{)}\OperatorTok{;}
\KeywordTok{end}

\CommentTok{\% True mean and variance of mixture}
\VariableTok{mu\_true} \OperatorTok{=} \FloatTok{0.8}\OperatorTok{*}\FloatTok{1} \OperatorTok{+} \FloatTok{0.2}\OperatorTok{*}\FloatTok{10}\OperatorTok{;}  \CommentTok{\% = 2.8}
\VariableTok{var\_true} \OperatorTok{=} \FloatTok{0.8}\OperatorTok{*}\NormalTok{(}\FloatTok{1}\OperatorTok{+}\FloatTok{1}\OperatorTok{\^{}}\FloatTok{2}\NormalTok{) }\OperatorTok{+} \FloatTok{0.2}\OperatorTok{*}\NormalTok{(}\FloatTok{100}\OperatorTok{+}\FloatTok{10}\OperatorTok{\^{}}\FloatTok{2}\NormalTok{) }\OperatorTok{{-}} \VariableTok{mu\_true}\OperatorTok{\^{}}\FloatTok{2}\OperatorTok{;}  \CommentTok{\% Second moment {-} mean\^{}2}

\CommentTok{\% Compute sample means}
\VariableTok{X\_bar} \OperatorTok{=} \VariableTok{mean}\NormalTok{(}\VariableTok{raw}\OperatorTok{,} \FloatTok{1}\NormalTok{)}\OperatorTok{;}

\VariableTok{figure}\OperatorTok{;}
\VariableTok{subplot}\NormalTok{(}\FloatTok{2}\OperatorTok{,}\FloatTok{1}\OperatorTok{,}\FloatTok{1}\NormalTok{)}\OperatorTok{;}
\VariableTok{histogram}\NormalTok{(}\VariableTok{raw}\NormalTok{(}\OperatorTok{:,}\FloatTok{1}\NormalTok{)}\OperatorTok{,} \FloatTok{100}\OperatorTok{,} \SpecialStringTok{\textquotesingle{}Normalization\textquotesingle{}}\OperatorTok{,} \SpecialStringTok{\textquotesingle{}pdf\textquotesingle{}}\NormalTok{)}\OperatorTok{;}
\VariableTok{title}\NormalTok{(}\SpecialStringTok{\textquotesingle{}Original Distribution (highly skewed mixture)\textquotesingle{}}\NormalTok{)}\OperatorTok{;}
\VariableTok{xlabel}\NormalTok{(}\SpecialStringTok{\textquotesingle{}X\textquotesingle{}}\NormalTok{)}\OperatorTok{;} \VariableTok{ylabel}\NormalTok{(}\SpecialStringTok{\textquotesingle{}PDF\textquotesingle{}}\NormalTok{)}\OperatorTok{;}

\VariableTok{subplot}\NormalTok{(}\FloatTok{2}\OperatorTok{,}\FloatTok{1}\OperatorTok{,}\FloatTok{2}\NormalTok{)}\OperatorTok{;}
\VariableTok{histogram}\NormalTok{(}\VariableTok{X\_bar}\OperatorTok{,} \FloatTok{100}\OperatorTok{,} \SpecialStringTok{\textquotesingle{}Normalization\textquotesingle{}}\OperatorTok{,} \SpecialStringTok{\textquotesingle{}pdf\textquotesingle{}}\NormalTok{)}\OperatorTok{;}
\VariableTok{hold} \VariableTok{on}\OperatorTok{;}
\VariableTok{x} \OperatorTok{=} \VariableTok{linspace}\NormalTok{(}\VariableTok{min}\NormalTok{(}\VariableTok{X\_bar}\NormalTok{)}\OperatorTok{,} \VariableTok{max}\NormalTok{(}\VariableTok{X\_bar}\NormalTok{)}\OperatorTok{,} \FloatTok{200}\NormalTok{)}\OperatorTok{;}
\VariableTok{plot}\NormalTok{(}\VariableTok{x}\OperatorTok{,} \VariableTok{normpdf}\NormalTok{(}\VariableTok{x}\OperatorTok{,} \VariableTok{mu\_true}\OperatorTok{,} \VariableTok{sqrt}\NormalTok{(}\VariableTok{var\_true}\OperatorTok{/}\VariableTok{n\_samples}\NormalTok{))}\OperatorTok{,} \SpecialStringTok{\textquotesingle{}r\textquotesingle{}}\OperatorTok{,} \SpecialStringTok{\textquotesingle{}LineWidth\textquotesingle{}}\OperatorTok{,} \FloatTok{2}\NormalTok{)}\OperatorTok{;}
\VariableTok{title}\NormalTok{(}\VariableTok{sprintf}\NormalTok{(}\SpecialStringTok{\textquotesingle{}Sample Mean (n=\%d) — Gaussian by CLT!\textquotesingle{}}\OperatorTok{,} \VariableTok{n\_samples}\NormalTok{))}\OperatorTok{;}
\VariableTok{xlabel}\NormalTok{(}\SpecialStringTok{\textquotesingle{}Sample Mean\textquotesingle{}}\NormalTok{)}\OperatorTok{;} \VariableTok{ylabel}\NormalTok{(}\SpecialStringTok{\textquotesingle{}PDF\textquotesingle{}}\NormalTok{)}\OperatorTok{;}
\VariableTok{legend}\NormalTok{(}\SpecialStringTok{\textquotesingle{}Empirical\textquotesingle{}}\OperatorTok{,} \SpecialStringTok{\textquotesingle{}Gaussian approximation\textquotesingle{}}\NormalTok{)}\OperatorTok{;}
\end{Highlighting}
\end{Shaded}

\PSsection{11.8}{تقریب‌زنی با استفاده از \lr{\mbox{CLT}}}

\PSsubsection{رویه کلی}

برای تقریب \lr{$P(S_{n} \le x)$} که در آن \lr{$S_{n} = X_{1} + \ldots + X_{n}$}:

\begin{enumerate}
\def\labelenumi{\arabic{enumi}.}
\tightlist
\item
  \lr{$\mu = E[X_{i}]$}، \lr{$\sigma^{2} = \operatorname{Var}(X_{i})$} را محاسبه کنید
\item
  \lr{$S_{n} \approx N(n\mu,\, n\sigma^{2})$}
\item
  \lr{$\displaystyle P(S_{n} \le x) \approx \Phi \left(\frac{x - n\mu}{\sigma \sqrt{n}}\right)$}
\end{enumerate}

\PSsubsection{تصحیح پیوستگی (برای متغیرهای تصادفی گسسته)}

هنگام تقریب‌زدن یک جمع گسسته با گاوسی:

\begin{itemize}
\tightlist
\item
  \lr{$\displaystyle P(X \le k) \approx \Phi \left(\frac{k + 0.5 - n\mu}{\sigma \sqrt{n}}\right)$}
\item
  \lr{$\displaystyle P(X = k) \approx \Phi \left(\frac{k + 0.5 - n\mu}{\sigma \sqrt{n}}\right) - \Phi \left(\frac{k - 0.5 - n\mu}{\sigma \sqrt{n}}\right)$}
\end{itemize}

\begin{PSbox}[unbreakable]{ex}{مثال: خطاهای بسته}

1000 بسته، احتمال خطای هر یک \PSnum{0.02}. \lr{$X$} = مجموع خطاها.\PSbr{}دقیق: \lr{$X \sim \mathrm{Binomial}(1000,\, 0.02)$}\PSbr{}تقریب \lr{\mbox{CLT}}: \lr{$X \approx N(20,\, 19.6)$}

\lr{$\displaystyle P(X > 30) \approx 1 - \Phi \left(\frac{30.5 - 20}{\sqrt{19.6}}\right) = 1 - \Phi (2.37) = 0.0089$}

\end{PSbox}

\PSsection{11.9}{مسایل تمرینی}

\begin{PSbox}[unbreakable]{prob}{مسیله 1}

یک کارخانه پیچ تولید می‌کند. طول \lr{$X$} \lr{\textasciitilde{}} توزیعی با \lr{$\mu = 10\,\mathrm{cm}$} و \lr{$\sigma = 0.2\,\mathrm{cm}$}. یک دسته 100 پیچی اندازه‌گیری می‌شود.

\begin{enumerate}
\def\labelenumi{(\PSalph{enumi})}
\tightlist
\item
  توزیع تقریبی \lr{$\bar{X}$} چیست؟ (ب) \lr{$P(\bar{X} > 10.05)$}؟ (ج) \lr{$P(9.96 < \bar{X} < 10.04)$}؟
\end{enumerate}

\PSsolution{} 

\begin{enumerate}
\def\labelenumi{(\PSalph{enumi})}
\tightlist
\item
  \lr{$\bar{X} \sim N\left(10,\, \frac{0.04}{100}\right) = N(10,\, 0.0004)$}، \lr{$\sigma_{\bar{X}} = 0.02$}
\item
  \lr{$P(\bar{X} > 10.05) = P(Z > 2.5) = 0.0062$}
\item
  \lr{$P(9.96 < \bar{X} < 10.04) = \Phi (2) - \Phi (-2) = 0.9545$}
\end{enumerate}

\end{PSbox}

\begin{PSbox}[unbreakable]{prob}{مسیله 2}

یک لینک مخابراتی 10000 بیت با \lr{$\mathrm{BER} = 0.001$} دارد. با استفاده از \lr{\mbox{CLT}} مقدار \lr{$P(\text{\rl{بیش از 15 خطا}})$} را تقریب بزنید.

\PSsolution{} \lr{\mbox{$X \sim \mathrm{Binomial}(10000,\, 0.001)$. $\mu = 10$}}، \lr{$\sigma^{2} = 9.99$}.\PSbr{}\lr{$\displaystyle P(X > 15) \approx 1 - \Phi \left(\frac{15.5-10}{\sqrt{9.99}}\right) = 1 - \Phi (1.74) = 0.0409$}

\end{PSbox}

\begin{PSbox}[unbreakable]{prob}{مسیله 3}

کد \lr{\mbox{MATLAB}} را برای بررسی \lr{\mbox{CLT}} برای توزیع کای‌دو(1) (بسیار چوله) بنویسید و همگرایی را برای \lr{$n = 2,\,5,\,10,\,50$} نشان دهید.

\PSsolution{} 

\tcbset{PScodemode/.style={unbreakable}}

\begin{Shaded}
\begin{Highlighting}[]
\VariableTok{N} \OperatorTok{=} \FloatTok{100000}\OperatorTok{;} \VariableTok{ns} \OperatorTok{=}\NormalTok{ [}\FloatTok{2} \FloatTok{5} \FloatTok{10} \FloatTok{50}\NormalTok{]}\OperatorTok{;}
\VariableTok{figure}\OperatorTok{;}
\KeywordTok{for} \VariableTok{i} \OperatorTok{=} \FloatTok{1}\OperatorTok{:}\FloatTok{4}
    \VariableTok{X\_bar} \OperatorTok{=} \VariableTok{mean}\NormalTok{(}\VariableTok{chi2rnd}\NormalTok{(}\FloatTok{1}\OperatorTok{,} \VariableTok{ns}\NormalTok{(}\VariableTok{i}\NormalTok{)}\OperatorTok{,} \VariableTok{N}\NormalTok{)}\OperatorTok{,} \FloatTok{1}\NormalTok{)}\OperatorTok{;}
    \VariableTok{subplot}\NormalTok{(}\FloatTok{2}\OperatorTok{,}\FloatTok{2}\OperatorTok{,}\VariableTok{i}\NormalTok{)}\OperatorTok{;}
    \VariableTok{histogram}\NormalTok{(}\VariableTok{X\_bar}\OperatorTok{,} \FloatTok{80}\OperatorTok{,} \SpecialStringTok{\textquotesingle{}Normalization\textquotesingle{}}\OperatorTok{,} \SpecialStringTok{\textquotesingle{}pdf\textquotesingle{}}\NormalTok{)}\OperatorTok{;} \VariableTok{hold} \VariableTok{on}\OperatorTok{;}
    \VariableTok{x} \OperatorTok{=} \VariableTok{linspace}\NormalTok{(}\FloatTok{0}\OperatorTok{,} \FloatTok{3}\OperatorTok{,} \FloatTok{200}\NormalTok{)}\OperatorTok{;}
    \VariableTok{plot}\NormalTok{(}\VariableTok{x}\OperatorTok{,} \VariableTok{normpdf}\NormalTok{(}\VariableTok{x}\OperatorTok{,} \FloatTok{1}\OperatorTok{,} \VariableTok{sqrt}\NormalTok{(}\FloatTok{2}\OperatorTok{/}\VariableTok{ns}\NormalTok{(}\VariableTok{i}\NormalTok{)))}\OperatorTok{,} \SpecialStringTok{\textquotesingle{}r\textquotesingle{}}\OperatorTok{,} \SpecialStringTok{\textquotesingle{}LineWidth\textquotesingle{}}\OperatorTok{,} \FloatTok{2}\NormalTok{)}\OperatorTok{;}
    \VariableTok{title}\NormalTok{(}\VariableTok{sprintf}\NormalTok{(}\SpecialStringTok{\textquotesingle{}n=\%d\textquotesingle{}}\OperatorTok{,} \VariableTok{ns}\NormalTok{(}\VariableTok{i}\NormalTok{)))}\OperatorTok{;}
\KeywordTok{end}
\VariableTok{sgtitle}\NormalTok{(}\SpecialStringTok{\textquotesingle{}CLT for χ²(1): Mean→Gaussian\textquotesingle{}}\NormalTok{)}\OperatorTok{;}
\end{Highlighting}
\end{Shaded}

\emph{ماژول بعدی: {ماژول 12 \lr{—} مقدمه‌ای بر آمار}}

\end{PSbox}
